% 2009 年浙江卷理科第 12 题：三视图（正视图、侧视图、俯视图）。
% 来源 PDF：普通高考/2009/2009浙江理.pdf，SHA-256 8b9a62d8af65c0049a04b2c462e5bc8dc66a464244ca2c2462c6cdfafc5c2d89。
% 原页：第 1 页第 2 栏中部，书页码 565。单位 cm。
% 重绘依据：正视图为 T 形，上横条宽 3（三段各 1）、高 1，竖柱宽 1、高 3；
%   侧视图为 3×4 矩形，距顶 1 处有一条横线；俯视图为 3×3 正方形，两条竖直虚线
%   为竖柱两侧棱的投影。按标注尺寸等比重绘（原页图线条未严格等比）。
\begin{tikzpicture}[semithick, every node/.style={inner sep=1.5pt, font=\small}]

  % ---------- 正视图 ----------
  \draw (0,4) -- (3,4) -- (3,3) -- (2,3) -- (2,0) -- (1,0) -- (1,3) -- (0,3) -- cycle;
  \foreach \a/\c in {0/0.5, 1/1.5, 2/2.5} {
    \draw[<->] (\a,4.28) -- (\a+1,4.28);
    \node[above] at (\c,4.28) {$1$};
  }
  \draw[<->] (3.30,3) -- (3.30,4);
  \node[right] at (3.30,3.5) {$1$};
  \draw[<->] (3.30,0) -- (3.30,3);
  \node[right] at (3.30,1.5) {$3$};
  \node at (1.5,-0.62) {正视图};

  % ---------- 侧视图 ----------
  \draw (4.90,4) -- (7.90,4) -- (7.90,0) -- (4.90,0) -- cycle;
  \draw (4.90,3) -- (7.90,3);
  \draw[<->] (4.90,4.28) -- (7.90,4.28);
  \node[above] at (6.40,4.28) {$3$};
  \node at (6.40,-0.62) {侧视图};

  % ---------- 俯视图 ----------
  \draw (0,-2) -- (3,-2) -- (3,-5) -- (0,-5) -- cycle;
  \draw[dashed] (1,-2) -- (1,-5);
  \draw[dashed] (2,-2) -- (2,-5);
  \draw[<->] (3.30,-2) -- (3.30,-5);
  \node[right] at (3.30,-3.5) {$3$};
  \node at (1.5,-5.62) {俯视图};
\end{tikzpicture}
